\section{Related work}
\label{sec:related}

\paragraph{Image Classification} is so core to computer vision that it is often used as a benchmark to measure progress in image understanding. 
Any progress usually translates to improvement in other related tasks such as detection or segmentation. 
Since 2012's AlexNet~\cite{Krizhevsky2012AlexNet}, convnets have dominated this benchmark and have become the de facto standard.
The evolution of the state of the art on the ImageNet dataset~\cite{Russakovsky2015ImageNet12} reflects the progress with convolutional neural network architectures and learning~\cite{Krizhevsky2012AlexNet,Simonyan2015VGG,tan2019efficientnet,Touvron2019FixRes,Touvron2020FixingTT,Xie2019SelftrainingWN}.

Despite several attempts to use transformers for image classification~\cite{chen2020generative}, until now their performance has been inferior to that of convnets.
Nevertheless hybrid architectures that combine convnets and transformers, including the self-attention mechanism, have recently exhibited competitive results in image classification~\cite{wu2020visual}, detection~\cite{carion2020end,Hu2018RelationNF}, video processing~\cite{Sun2019VideoBERTAJ,Wang2018NonlocalNN}, unsupervised object discovery~\cite{Locatello2020ObjectCentricLW}, and unified text-vision tasks~\cite{Chen2020UNITERUI,li2019visualbert,Lu2019ViLBERTPT}.



Recently Vision transformers (ViT)~\cite{dosovitskiy2020image} closed the gap with the state of the art on ImageNet, without using any convolution.
This performance is remarkable since convnet methods for image classification have benefited from years of tuning and optimization~\cite{he2019bag,pytorchmodels}.
Nevertheless, according to this study~\cite{dosovitskiy2020image}, a pre-training phase on a large volume of curated data is required for the learned transformer to be effective. 
% 
In our paper we achieve a strong performance without requiring a large training dataset, i.e., with Imagenet1k only. 
%
\paragraph{The Transformer architecture,}
%
introduced by Vaswani et al.~\cite{Vaswani2017AttentionIA} for machine translation are currently the reference model for all natural language processing (NLP) tasks.
Many improvements of convnets for image classification are inspired by transformers.
For example, Squeeze and Excitation~\cite{Hu2017SENet}, Selective Kernel~\cite{Li2019SelectiveKN} and Split-Attention Networks~\cite{zhang2020resnest} exploit mechanism akin to transformers self-attention (SA) mechanism.

\paragraph{Knowledge Distillation} (KD), introduced by Hinton et al.~\cite{Hinton2015DistillingTK}, refers to the training paradigm in which a \emph{student} model leverages ``soft'' labels coming from a strong \emph{teacher} network. 
This is the output vector of the teacher's softmax function rather than just the maximum of scores, wich gives a ``hard'' label.
Such a training improves the performance of the student model (alternatively, it can be regarded as a form of compression of the teacher model into a smaller one -- the student).
On the one hand the teacher's soft labels will have a similar effect to labels smoothing~\cite{Yuan2019RevisitKD}. %
On the other hand as shown by Wei et al.~\cite{Wei2020CircumventingOO} the teacher's supervision takes into account the effects of the data augmentation, which sometimes causes a misalignment between the real label and the image.
For example, let us consider image with a ``cat'' label that represents a large landscape and a small cat in a corner. 
If the cat is no longer on the crop of the data augmentation it implicitly changes the label of the image. 
KD can transfer inductive biases~\cite{abnar2020transferring}  in a soft way in a student model using a teacher model where they would be incorporated in a hard way. 
For example, it may be useful to induce biases due to convolutions in a transformer model by using a convolutional model as teacher. 
In our paper we study the distillation of a transformer student by either a convnet or a transformer teacher. We introduce a new distillation procedure specific to transformers and show its superiority. 
